% Chapter 6 — Data Architecture (Phase 2). SCAFFOLD: write the prose into each section.
% Depth cap: section -> subsection only (no subsubsection).

\section{Architecture Overview}
\label{sec:arch:overview}
% Write: the production data architecture — dbt (transforms) + Dagster (orchestration) over
% BigQuery, with DuckDB for local sampling. Storage/compute separation, with dbt as the
% portability seam. Include one architecture diagram (draw.io).

\section{Medallion Warehouse}
\label{sec:arch:medallion}
% Write: Bronze sources (parquet on GCS, read as BigQuery external tables) -> Silver Kimball star
% (dim_* tables + fact_* views) -> Gold aggregates. Incremental, partitioned by day, clustered by
% route_no / fare_class. The three Silver integrity fixes. Integrity as dbt tests plus the
% relative lossless test (Gold boardings == Bronze ticket count).

\section{Ingestion and Orchestration}
\label{sec:arch:ingestion}
% Write: micro-batch ingestion (daily/hourly exports -> GCS) -> Dagster sensor -> incremental
% dbt run. The filterable Gold cube (date x route x day-type x fare-class x hour) the filter-bar
% endpoints read. No streaming-analytics engine: the product is batch aggregate analytics.

\section{Engine Portability and Scale}
\label{sec:arch:scale}
% Write: BigQuery (bq_dev sampled / bq_prod full) as the primary engine; DuckDB as a local
% sampling target. Dynamic, continuously-growing data; relative invariants. The scale rationale
% (millions of trips/tickets per day).
